% Options for packages loaded elsewhere
% Options for packages loaded elsewhere
\PassOptionsToPackage{unicode}{hyperref}
\PassOptionsToPackage{hyphens}{url}
\PassOptionsToPackage{dvipsnames,svgnames,x11names}{xcolor}
%
\documentclass[
  japanese,
  letterpaper,
  DIV=11,
  numbers=noendperiod]{scrartcl}
\usepackage{xcolor}
\usepackage{amsmath,amssymb}
\setcounter{secnumdepth}{-\maxdimen} % remove section numbering
\usepackage{iftex}
\ifPDFTeX
  \usepackage[T1]{fontenc}
  \usepackage[utf8]{inputenc}
  \usepackage{textcomp} % provide euro and other symbols
\else % if luatex or xetex
  \usepackage{unicode-math} % this also loads fontspec
  \defaultfontfeatures{Scale=MatchLowercase}
  \defaultfontfeatures[\rmfamily]{Ligatures=TeX,Scale=1}
\fi
\usepackage{lmodern}
\ifPDFTeX\else
  % xetex/luatex font selection
\fi
% Use upquote if available, for straight quotes in verbatim environments
\IfFileExists{upquote.sty}{\usepackage{upquote}}{}
\IfFileExists{microtype.sty}{% use microtype if available
  \usepackage[]{microtype}
  \UseMicrotypeSet[protrusion]{basicmath} % disable protrusion for tt fonts
}{}
\makeatletter
\@ifundefined{KOMAClassName}{% if non-KOMA class
  \IfFileExists{parskip.sty}{%
    \usepackage{parskip}
  }{% else
    \setlength{\parindent}{0pt}
    \setlength{\parskip}{6pt plus 2pt minus 1pt}}
}{% if KOMA class
  \KOMAoptions{parskip=half}}
\makeatother
% Make \paragraph and \subparagraph free-standing
\makeatletter
\ifx\paragraph\undefined\else
  \let\oldparagraph\paragraph
  \renewcommand{\paragraph}{
    \@ifstar
      \xxxParagraphStar
      \xxxParagraphNoStar
  }
  \newcommand{\xxxParagraphStar}[1]{\oldparagraph*{#1}\mbox{}}
  \newcommand{\xxxParagraphNoStar}[1]{\oldparagraph{#1}\mbox{}}
\fi
\ifx\subparagraph\undefined\else
  \let\oldsubparagraph\subparagraph
  \renewcommand{\subparagraph}{
    \@ifstar
      \xxxSubParagraphStar
      \xxxSubParagraphNoStar
  }
  \newcommand{\xxxSubParagraphStar}[1]{\oldsubparagraph*{#1}\mbox{}}
  \newcommand{\xxxSubParagraphNoStar}[1]{\oldsubparagraph{#1}\mbox{}}
\fi
\makeatother


\usepackage{longtable,booktabs,array}
\usepackage{calc} % for calculating minipage widths
% Correct order of tables after \paragraph or \subparagraph
\usepackage{etoolbox}
\makeatletter
\patchcmd\longtable{\par}{\if@noskipsec\mbox{}\fi\par}{}{}
\makeatother
% Allow footnotes in longtable head/foot
\IfFileExists{footnotehyper.sty}{\usepackage{footnotehyper}}{\usepackage{footnote}}
\makesavenoteenv{longtable}
\usepackage{graphicx}
\makeatletter
\newsavebox\pandoc@box
\newcommand*\pandocbounded[1]{% scales image to fit in text height/width
  \sbox\pandoc@box{#1}%
  \Gscale@div\@tempa{\textheight}{\dimexpr\ht\pandoc@box+\dp\pandoc@box\relax}%
  \Gscale@div\@tempb{\linewidth}{\wd\pandoc@box}%
  \ifdim\@tempb\p@<\@tempa\p@\let\@tempa\@tempb\fi% select the smaller of both
  \ifdim\@tempa\p@<\p@\scalebox{\@tempa}{\usebox\pandoc@box}%
  \else\usebox{\pandoc@box}%
  \fi%
}
% Set default figure placement to htbp
\def\fps@figure{htbp}
\makeatother



\ifLuaTeX
\usepackage[bidi=basic,provide=*]{babel}
\else
\usepackage[bidi=default,provide=*]{babel}
\fi
% get rid of language-specific shorthands (see #6817):
\let\LanguageShortHands\languageshorthands
\def\languageshorthands#1{}


\setlength{\emergencystretch}{3em} % prevent overfull lines

\providecommand{\tightlist}{%
  \setlength{\itemsep}{0pt}\setlength{\parskip}{0pt}}



 


\usepackage{booktabs}
\usepackage{longtable}
\usepackage{array}
\usepackage{multirow}
\usepackage{wrapfig}
\usepackage{float}
\usepackage{colortbl}
\usepackage{pdflscape}
\usepackage{tabu}
\usepackage{threeparttable}
\usepackage{threeparttablex}
\usepackage[normalem]{ulem}
\usepackage{makecell}
\usepackage{xcolor}
\usepackage{fontspec}
\usepackage{xeCJK}
\usepackage{booktabs}
\usepackage{caption}
\captionsetup[table]{name=Table,labelsep=period,skip=6pt}
\usepackage{longtable}
\captionsetup[longtable]{skip=0pt}
\usepackage{xcolor}
\KOMAoption{captions}{tableheading}
\makeatletter
\@ifpackageloaded{caption}{}{\usepackage{caption}}
\AtBeginDocument{%
\ifdefined\contentsname
  \renewcommand*\contentsname{目次}
\else
  \newcommand\contentsname{目次}
\fi
\ifdefined\listfigurename
  \renewcommand*\listfigurename{図一覧}
\else
  \newcommand\listfigurename{図一覧}
\fi
\ifdefined\listtablename
  \renewcommand*\listtablename{表一覧}
\else
  \newcommand\listtablename{表一覧}
\fi
\ifdefined\figurename
  \renewcommand*\figurename{図}
\else
  \newcommand\figurename{図}
\fi
\ifdefined\tablename
  \renewcommand*\tablename{表}
\else
  \newcommand\tablename{表}
\fi
}
\@ifpackageloaded{float}{}{\usepackage{float}}
\floatstyle{ruled}
\@ifundefined{c@chapter}{\newfloat{codelisting}{h}{lop}}{\newfloat{codelisting}{h}{lop}[chapter]}
\floatname{codelisting}{コード}
\newcommand*\listoflistings{\listof{codelisting}{コード一覧}}
\makeatother
\makeatletter
\makeatother
\makeatletter
\@ifpackageloaded{caption}{}{\usepackage{caption}}
\@ifpackageloaded{subcaption}{}{\usepackage{subcaption}}
\makeatother
\usepackage{bookmark}
\IfFileExists{xurl.sty}{\usepackage{xurl}}{} % add URL line breaks if available
\urlstyle{same}
\hypersetup{
  pdftitle={kableExtra LaTeX チートシート},
  pdfauthor={Generated by ChatGPT},
  pdflang={ja},
  colorlinks=true,
  linkcolor={blue},
  filecolor={Maroon},
  citecolor={Blue},
  urlcolor={Blue},
  pdfcreator={LaTeX via pandoc}}


\title{kableExtra LaTeX チートシート}
\usepackage{etoolbox}
\makeatletter
\providecommand{\subtitle}[1]{% add subtitle to \maketitle
  \apptocmd{\@title}{\par {\large #1 \par}}{}{}
}
\makeatother
\subtitle{論文PDF向け・使い方とコード例}
\author{Generated by ChatGPT}
\date{2025-11-12}
\begin{document}
\maketitle


\section{使い方の前提}\label{ux4f7fux3044ux65b9ux306eux524dux63d0}

\begin{itemize}
\tightlist
\item
  本書は \textbf{LaTeX（PDF）出力}に特化した \emph{kableExtra}
  の実用チートシートです。\\
\item
  Quarto(.qmd) を \textbf{Render}（Knitではなく）してください。\\
\item
  事前に \texttt{library(kableExtra)}
  などを読み込めば、そのまま実行できます。
\end{itemize}

\section{1. 最小例（booktabs \& 段組み \&
脚注）}\label{ux6700ux5c0fux4f8bbooktabs-ux6bb5ux7d44ux307f-ux811aux6ce8}

\begin{table}

\caption{\label{tbl-minimal}最小レシピ：段組みと脚注}

\centering{

\centering
\begin{threeparttable}
\begin{tabular}[t]{lrrr}
\toprule
\multicolumn{1}{c}{ } & \multicolumn{3}{c}{Summary} \\
\cmidrule(l{3pt}r{3pt}){2-4}
Group & $n$ & Mean & SD\\
\midrule
Control & 30 & 12.34 & 2.11\\
Treatment A & 28 & 13.21 & 1.87\\
Treatment B & 27 & 15.01 & 2.53\\
\bottomrule
\end{tabular}
\begin{tablenotes}
\item \textit{Note: } 
\item Values are mean (SD).
\end{tablenotes}
\end{threeparttable}

}

\end{table}%

本文からは Table~\ref{tbl-minimal} のように参照できます。

\section{2.
主な関数と重要引数（LaTeX版）}\label{ux4e3bux306aux95a2ux6570ux3068ux91cdux8981ux5f15ux6570latexux7248}

\begin{table}[!h]
\centering
\resizebox{\ifdim\width>\linewidth\linewidth\else\width\fi}{!}{
\begin{tabular}[t]{llll}
\toprule
Function & Key arguments & Purpose & Notes\\
\midrule
kbl() & format, booktabs, longtable, digits, col.names, row.names, align, caption, label, escape, format.args, valign, position & 基礎テーブル生成（knitr::kableの拡張） & booktabs=TRUE で綺麗な罫線。escape=FALSE で数式OK。\\
kable\_styling() & latex\_options (hold\_position, striped, scale\_down, repeat\_header), font\_size, position, full\_width & 全体スタイルや配置を整える & repeat\_header は longtable=TRUE と併用。\\
add\_header\_above() & header(named vector), bold, italic, align, line(LaTeX), escape & 上段の列見出し（段組み）を追加 & 合計列数=表の列数に一致。LaTeXマクロは escape=FALSE。\\
footnote() & general, number, symbol, alphabet, threeparttable & 三部形式の脚注を追加 & threeparttable を自動利用。\\
column\_spec() & column, width('2cm'), bold, italic, color, background, latex\_column\_spec & 列幅や装飾を列単位で制御 & xcolor が自動で読み込まれる。\\
\addlinespace
row\_spec() & row(0=header), bold, color, background, hline\_after & 行単位の装飾や罫線制御 & ヘッダ行は row=0。部分罫線は `\textbackslash{}cmidrule` を併用可。\\
group\_rows() & group\_label, start\_row, end\_row, escape & 連続行を帯でグループ化 & 行番号はソート後のインデックスで指定。\\
collapse\_rows() & columns, latex\_hline, valign, target & 同値セルの縦結合（連続部分） & 長表との相性は要検証。まずは短表で確認。\\
landscape() & margin & 横向きページで出力 & pdflscape 依存。横長表に便利。\\
save\_kable() & file, latex\_header\_includes & テーブルをファイル保存 & PDFは通常 knit 経由推奨。\\
\bottomrule
\end{tabular}}
\end{table}

\section{3. 列幅・色・行線}\label{ux5217ux5e45ux8272ux884cux7dda}

\begin{table}

\caption{\label{tbl-styling}列幅・色・行線の例}

\centering{

\centering
\begin{tabular}[t]{>{\raggedright\arraybackslash}p{3cm}>{\raggedleft\arraybackslash}p{1.6cm}>{\raggedleft\arraybackslash}p{1.6cm}>{\raggedleft\arraybackslash}p{1.6cm}r}
\toprule
\textbf{ } & \textbf{mpg} & \textbf{cyl} & \textbf{disp} & \textbf{hp}\\
\midrule
\textbf{Mazda RX4} & 21.0 & \textcolor{blue}{6} & \cellcolor[HTML]{EFEFEF}{160} & 110\\
\textbf{Mazda RX4 Wag} & 21.0 & \textcolor{blue}{6} & \cellcolor[HTML]{EFEFEF}{160} & 110\\
\textbf{Datsun 710} & 22.8 & \textcolor{blue}{4} & \cellcolor[HTML]{EFEFEF}{108} & 93\\
\midrule
\textbf{Hornet 4 Drive} & 21.4 & \textcolor{blue}{6} & \cellcolor[HTML]{EFEFEF}{258} & 110\\
\textbf{Hornet Sportabout} & 18.7 & \textcolor{blue}{8} & \cellcolor[HTML]{EFEFEF}{360} & 175\\
\addlinespace
\textbf{Valiant} & 18.1 & \textcolor{blue}{6} & \cellcolor[HTML]{EFEFEF}{225} & 105\\
\bottomrule
\end{tabular}

}

\end{table}%

\section{4.
longtable（複数ページ）}\label{longtableux8907ux6570ux30daux30fcux30b8}

\begin{longtable}[t]{rrr}

\caption{\label{tbl-longtable}longtable
の基本（ページ跨ぎでヘッダ繰り返し）}

\tabularnewline

\toprule
A & B & C\\
\midrule
\endfirsthead
\multicolumn{3}{@{}l}{\textit{(continued)}}\\
\toprule
A & B & C\\
\midrule
\endhead

\endfoot
\bottomrule
\endlastfoot
1 & 0.7690785 & -0.6401028\\
2 & 0.0945735 & 0.5835009\\
3 & 0.5400228 & 0.1430305\\
4 & 0.4757187 & 3.4147689\\
5 & 1.0108121 & 0.3688312\\
\addlinespace
6 & 0.0148302 & -2.2259096\\
7 & -0.9814969 & -0.4182953\\
8 & 1.9256670 & -2.0654373\\
9 & -1.2315970 & -0.4550973\\
10 & 1.4865365 & -0.5628451\\
\addlinespace
11 & -0.0116775 & 2.0868347\\
12 & 0.8279345 & -1.2591847\\
13 & 0.8243959 & -1.1612956\\
14 & 0.7538591 & -0.7651515\\
15 & -0.8806291 & 0.7244821\\
\addlinespace
16 & 0.7981530 & -1.0306616\\
17 & 0.1783338 & 0.7096168\\
18 & 0.1749529 & -0.0239833\\
19 & -0.2828129 & -1.1746491\\
20 & -0.2858078 & -0.2647614\\
\addlinespace
21 & 2.3170450 & 0.7351135\\
22 & 2.3415760 & 1.2355407\\
23 & -0.0618315 & 0.5986492\\
24 & 0.0031679 & 0.4960013\\
25 & 0.1366049 & 0.7621820\\
\addlinespace
26 & 1.1049060 & 0.0217209\\
27 & -0.4481680 & -0.3821367\\
28 & 0.7260534 & 1.4248137\\
29 & -0.8557726 & 0.3245495\\
30 & -0.2344179 & 0.6680187\\
\addlinespace
31 & 0.1835784 & 2.1060576\\
32 & 0.7535756 & -0.2077767\\
33 & 1.1100167 & 1.3210180\\
34 & 0.4178697 & -0.0707841\\
35 & 0.3547993 & -0.6537265\\
\addlinespace
36 & 0.1216973 & 2.4652743\\
37 & 0.6534702 & -0.9785047\\
38 & 1.2020585 & -0.2066036\\
39 & -0.5028535 & -0.2499698\\
40 & -0.1873698 & -2.1378912\\
\addlinespace
41 & 0.3461981 & -0.5982500\\
42 & -0.5702242 & 0.4694985\\
43 & -1.2499020 & -1.5245718\\
44 & 0.8181522 & -0.5508928\\
45 & -0.5700436 & -0.4612399\\
\addlinespace
46 & -0.5949977 & -1.5080083\\
47 & -0.2890411 & 1.8525430\\
48 & -0.3918290 & -0.6940291\\
49 & 1.3459708 & 0.8584311\\
50 & 1.1371514 & 0.2467644\\
\addlinespace
51 & -0.7131617 & -0.6465214\\
52 & 2.3062345 & -0.5896288\\
53 & -0.1438764 & 0.0292732\\
54 & -1.4236218 & 0.8999554\\
55 & -0.1588518 & 0.8814837\\
\addlinespace
56 & -0.9334786 & 1.6937897\\
57 & 0.0831327 & -1.8017939\\
58 & 0.6267968 & 1.2657011\\
59 & -0.7306676 & -0.2635127\\
60 & 1.1688684 & 1.1297121\\
\addlinespace
61 & -0.8241676 & -1.1200201\\
62 & 0.1504654 & -1.6573005\\
63 & -0.5042832 & -1.0865505\\
64 & -0.7397871 & 0.2573165\\
65 & -0.2908912 & 0.9015591\\
\addlinespace
66 & -0.4675384 & -0.2163925\\
67 & -0.4098618 & 0.2522726\\
68 & 0.7673985 & -0.2334731\\
69 & -0.9064484 & 2.1804813\\
70 & -0.5175915 & 0.3183159\\
\addlinespace
71 & 0.5412852 & 0.3838292\\
72 & -0.0542327 & -0.5924157\\
73 & 0.9548811 & 0.1703646\\
74 & -0.9551514 & 0.6648930\\
75 & -0.1330722 & -1.2799966\\
\addlinespace
76 & -0.6970350 & -0.5541463\\
77 & 1.3790795 & 0.5973675\\
78 & 0.7601221 & -0.0606880\\
79 & 0.4845187 & 0.0501421\\
80 & -0.8847302 & 0.8098824\\
\addlinespace
81 & -0.3617090 & -0.7845303\\
82 & 1.1937214 & -0.8079587\\
83 & -0.3365179 & 1.1019131\\
84 & -1.4994061 & 0.4526525\\
85 & 0.6883861 & -0.8778149\\
\addlinespace
86 & -1.2565035 & 0.5404907\\
87 & -0.4622839 & -0.6736491\\
88 & -1.6957714 & 0.3246053\\
89 & 1.6033491 & 1.3955095\\
90 & -0.0773196 & 0.8537696\\
\addlinespace
91 & -0.2660694 & -0.5392963\\
92 & 0.5313561 & -0.0670908\\
93 & 0.4540572 & 0.2906668\\
94 & 0.5884208 & -0.6347528\\
95 & 1.5466024 & -0.0920957\\
\addlinespace
96 & -1.4052828 & -0.1068115\\
97 & -1.0773535 & -1.0822714\\
98 & 1.0933272 & -0.2552253\\
99 & 2.2214450 & 0.5601084\\
100 & 0.2772038 & -0.4288075\\
\addlinespace
101 & 0.5168891 & -0.6498781\\
102 & 0.0815673 & 1.0619248\\
103 & 0.2634253 & 0.0506282\\
104 & -1.6147579 & 1.3448236\\
105 & -1.0716271 & -0.0201085\\
\addlinespace
106 & 0.8478413 & -0.5665250\\
107 & -1.0400572 & -0.3220919\\
108 & -0.6635421 & -0.2545636\\
109 & 1.2354759 & -1.0787687\\
110 & 1.3740948 & 0.9567265\\
\addlinespace
111 & 0.3327897 & 0.4498769\\
112 & -0.2569577 & 0.0534348\\
113 & 0.6725709 & -1.7022170\\
114 & 1.1167893 & -1.6115824\\
115 & -0.5851974 & 0.6238670\\
\addlinespace
116 & -0.0547761 & 0.0375677\\
117 & 0.1215148 & -1.9340150\\
118 & 0.9837539 & -0.4911344\\
119 & 1.1017398 & 0.7534654\\
120 & -0.0329875 & 0.2393494\\*

\end{longtable}

\section{5. よくある落とし穴 \&
回避}\label{ux3088ux304fux3042ux308bux843dux3068ux3057ux7a74-ux56deux907f}

\begin{itemize}
\tightlist
\item
  \textbf{キャプションが出ない}：Quarto の
  \texttt{\#\textbar{}\ tbl-cap:} は \textbf{チャンク先頭}に置く。\\
\item
  \textbf{LaTeXコードが文字のまま}：\texttt{escape=FALSE} を付ける。\\
\item
  \textbf{表が本文から離れる}：\texttt{kable\_styling(latex\_options=\textquotesingle{}hold\_position\textquotesingle{})}。\\
\item
  \textbf{``Table → Tab.''などの語を変えたい}：YAMLの
  \texttt{include-in-header} で
  \texttt{\textbackslash{}captionsetup{[}table{]}\{name=Tab.,labelsep=period\}}。\\
\item
  \textbf{キャプションと表の間隔}：\texttt{\textbackslash{}captionsetup{[}table{]}\{skip=6pt\}}。\\
\item
  \textbf{行の部分罫線}：\texttt{row\_spec(0,\ extra\_latex\_after=\textquotesingle{}\textbackslash{}\textbackslash{}cmidrule(lr)\{2-4\}\textquotesingle{})}（版によっては
  add\_header\_above 直後に
  \texttt{\textbackslash{}\textbackslash{}cmidrule} を手書き）。
\end{itemize}

\section{6.
付録：最小テンプレ（コピペ用）}\label{ux4ed8ux9332ux6700ux5c0fux30c6ux30f3ux30d7ux30ecux30b3ux30d4ux30daux7528}

\begin{table}

\caption{\label{tbl-template}最小テンプレ}

\centering{

\centering
\begin{tabular}[t]{lrrr}
\toprule
\multicolumn{1}{c}{ } & \multicolumn{3}{c}{Summary} \\
\cmidrule(l{3pt}r{3pt}){2-4}
Group & $n$ & Mean & SD\\
\midrule
A & 10 & 1.23 & 0.11\\
B & 12 & 4.56 & 0.22\\
\bottomrule
\end{tabular}

}

\end{table}%




\end{document}
